\chapter{Defeaters, Limitations, and Intellectual Risks}

The framework is designed to become narrower when evidence fails rather than to preserve its preferred interpretation at all costs. Several limitations therefore belong to the architecture rather than to an apologetic final paragraph.

The first risk is trivial persistence. Any conserved or slowly varying quantity can be used to construct a form map under which states appear persistent. The response is that $\Phi$ must be tied to a scientifically justified identity criterion, causal organization, invariant structure, or interventionally relevant observable. Persistence is not established by finding something that remains constant; the preserved relation must be the relation whose survival matters to the system being studied.

The second risk is hidden dualism. Writing $P_{\mathrm{phys}}\leftarrow\mathscr R\rightarrow E_{\mathrm{exp}}$ can look like property dualism with an additional latent variable. The dual-aspect proposal avoids trivial dualism only if $\mathscr R$ has independent mathematical structure, both projections are constrained by that structure, and the joint model produces consequences not obtainable merely by storing two unrelated coordinate sets. Without those conditions the notation contributes no explanation.

The third risk is an occupancy escape hatch. If every positive residual is treated as evidence for a complementary aspect while every negative result is dismissed as incomplete physics, the program cannot lose. The primary defeater is therefore physical absorption: if systematic enrichment of $P_{\mathrm{phys}}^{(k)}$ drives $\Delta_\varphi(k)$ toward zero robustly across enrichment orders, contexts, and interventions, the tested domain offers no empirical support for irreducible complementarity. A second defeater is latent instability: if $Z_\varphi$ changes substantially when one measurement channel is removed, it lacks independent phenomenological anchoring. A third is failure of transfer: if each context requires a new bridge map, no invariant cross-aspect relation has been discovered. A fourth is capacity equivalence: if a complexity-matched physical latent model predicts as well as the relational model, the latter gains no empirical advantage. A fifth is transformation nonspecificity: if unitary or generic nonlinear bridges fit as well as the antiunitary model, the antiunitary subhypothesis fails.

The fourth risk is Platonic mystique. Finite rotation groups and spherical harmonics belong to standard representation theory and equivariant pattern formation, not to a privileged sacred geometry. Their role here is technical: they provide exact examples of invariant form, admissibility, multiplicity, and nonlinear selection. The book should never use the amount of Platonic mathematics as evidence that reality is ontologically Platonic.

The fifth risk is conjugation mystique. Complex conjugation provides a precise mathematical involution with density preservation and current reversal in the scalar wave setting. It is not spirit, an afterlife state, a hidden observer, or consciousness. Any experiential application requires an independently justified complex phenomenological geometry and a falsifiable transformation test.

The sixth risk is dimension counting as ontology. The rank of projections can show local complementarity but cannot show interiority. Likewise, a kernel or flat direction can show a failure of a chosen description to distinguish states but cannot show that an unmeasured dimension is experiential. The $\ell=5$ flat directions are explicitly physical degeneracies of a truncated invariant energy. The new sextic dimension count strengthens the algebraic resolution result, but four primitive sextic directions remain mathematical structure, not evidence for a metaphysical dimension.

The seventh risk is confusing ambient and intrinsic curvature. A tangent vector to a constrained manifold does not imply that the straight line in that tangent direction remains on the manifold. For $J_{4,8}$, the normalized ambient-Hessian contribution is $1094/2145$, while the true constrained shape Hessian is $336/715$ after the second-order correction required by both $\|A\|=1$ and $B_2(A)=0$. Any future constrained stability calculation must include the appropriate second-fundamental-form term rather than treating a restricted ambient Hessian as intrinsic by default.

The eighth risk is overclaiming physical nonclosure. A measured physical projection may fail to close because of ordinary hidden physical variables, memory, insufficient state reconstruction, or noise. Approximate closure under a declared horizon and error tolerance is therefore more relevant experimentally than an idealized exact closure claim.

The ninth risk is underdetermination of causal or ontological direction. The most likely outcome of a mature ontology experiment may be that several directed models remain empirically equivalent. Observed relation does not by itself determine direction of predication. Physical intervention constrains viable models but does not automatically establish physical ontological priority, and experiential immediacy does not establish the reverse arrow. The project therefore adopts smaller success criteria: improved mathematical constraints, eliminated bridge theories, better phenomenological measurement, transferable causal laws, and clearer empirical boundaries all count as progress even if final ontological interpretation remains open.

The tenth limitation is that the PDE-specific sixth-order calculation for $\ell=5$ remains incomplete. What is now closed is stronger than a dimension count: the primitive sextic quotient has dimension four, the explicit $J_{4,8}$ witness has been independently reconstructed, and its positive-definite intrinsic shape Hessian has been independently reproduced exactly. What remains open is whether the actual PDE's fifth-order amplitude reduction produces a nonzero resolving component. The first sharp dynamical target is the restricted matrix $H_{6,\mathrm{shape}}^{\mathrm{PDE}}$, not an assumed full $H_6$.

The eleventh risk is recursive-geometry overreach. Appendix G uses Apollonian packing and graph-within-graph constructions as expository models of recursive refinement. Descartes' Circle Theorem establishes that a finite local rule can generate indefinitely fine recursive structure in that mathematical construction. It does not establish that physical reality is recursively packed, that the defined state $\mathcal X=(\{C_i\},T,\mathcal R,\Phi)$ is ontologically fundamental, or that nested geometry supports any conclusion about consciousness. The appendix is deliberately separated from the theorem-bearing $\ell=5$ audit for this reason.

Finally, the present choice of relational state and invariant mathematics is itself provisional. Other mathematical languages may outperform it for some questions. The framework must therefore accept replacement when another formalism preserves the relevant structure more economically or exposes distinctions the current one cannot express. That willingness is not a weakness of the relational program; it is the direct consequence of taking descriptive adequacy seriously.
